\begin{lemma}[Well-founded prefix tree]
For every front $F$, the proper-extension relation on its prefix tree
$T(F)$ is well-founded.  Equivalently, $T(F)$ has no infinite branch, so
there is no strictly increasing sequence $x(0)<x(1)<\cdots$ such that every
finite initial set $\{x(0),\ldots,x(n-1)\}$ belongs to $T(F)$.
\end{lemma}
\begin{proof}
Suppose an infinite branch existed and let $X'$ be its infinite range.  The
strict increase makes its finite nodes exactly the initial segments of
$X'$.  By the definition of \noderef{PrefixTree}, every such node is an
initial segment of a member of $F$, so $X'\subseteq\mathsf{FrontBase}(F)$.

If $F$ is trivial, its prefix tree contains only $\emptyset$, contradicting
the branch at length one.  Otherwise the base clause of \noderef{Front} gives
$X'\subseteq X$.  Density in \noderef{Front} supplies $s\in F$ with
$s\segs X'$.  Let $u$ be the next finite node on the branch after the node
$s$; then $s$ is a proper initial segment of $u$.  Since $u\in T(F)$,
\noderef{PrefixTree} supplies $t\in F$ with $u$ an initial segment of $t$.
Thus $s$ is an initial segment of $t$.  Prefix-freeness in
\noderef{Front}, together with $s,t\in F$, gives $s=t$, while
$s\segs u\segm t$ gives $s\ne t$, a contradiction.  Hence there is no
infinite branch, and the proper-extension relation is well-founded.
\end{proof}
